\documentclass[11pt]{article}

\usepackage[margin=25mm]{geometry}

\usepackage{amsmath,amssymb,fontspec,unicode-math}

\IfFontExistsTF{DejaVu Serif}{\setmainfont{DejaVu Serif}}{\setmainfont{Latin Modern Roman}}

\IfFontExistsTF{Latin Modern Math}{\setmathfont{Latin Modern Math}}{\setmathfont{DejaVu Math TeX Gyre}}

\usepackage{xurl,hyperref}

\setlength{\parindent}{0pt}\setlength{\parskip}{6pt}

\begin{document}

\title{Policy Compatible Skill Incremental Learning via Lazy Learning Interface}

\author{ELVIA preliminary AI analysis}\date{21 October 2026}\maketitle

Preliminary assessment: Minor.\par

Independent scientific review remains pending. Author adoption has not been requested.

The latest paper’s unconditional fallback union can select an already rejected skill despite an available goal-valid skill. The official supplement and current author implementation already use a conditional candidate branch that avoids this counterexample.

\section*{Review scope}

Claim-focused reading of the original paper and relevant author-released source, with direct mathematical and static-source checks. No author program, historical verifier, model training or published benchmark was executed.

\section*{Problem 1: An unconditional fallback can defeat goal validation}

Source: Section 4.3, Eq.(8), p.5; Appendix A.2, Algorithm 2; Appendix A.3, Figure 8; lazySI.py, forward, lines 1111–1154

Assessment: supported; impact: minor

\subsection*{Source claim (paraphrase)}

When the proposed skill fails goal validation, the printed rule ranks state proximity over \(A(g)\cup\{z_h\}\), where \(A(g)\) contains goal-valid skills. The following prose says the rejected proposal \(z_h\) is only a fallback when no valid candidate exists.

\subsection*{Analysis}

Take two skills, proposal \(z_h=1\), state \(s=0\), and goal \(g=10\). Give the state prototypes means \(0,1\), the goal prototypes means \(0,10\), and every Gaussian variance \(1\). With threshold \(\sqrt{\chi^2_{1,0.99}}\approx2.575829\), skill \(1\) fails goal validation and skill \(2\) passes. The unconditional union nevertheless ranks skill \(1\) first because its state distance is \(0\), versus \(1\). Four duplicate prototypes satisfy the stated prototype count; small perturbations preserve the strict inequalities. The same construction embeds in \(60\) dimensions with identity covariance, whose threshold is approximately \(9.401<10\). In contrast, the author implementation collects valid candidates and introduces its nearest-goal fallback only when that collection is empty, so it avoids this example.

\subsection*{Finding and scope}

The current printed candidate formula does not enforce its own stated fallback semantics. This is a paper/implementation specification inconsistency, with no demonstrated current-code failure or reward reduction.

\subsection*{Proposed correction}

Make the rule conditional: retain \(z_h\) if it is valid; otherwise rank over \(A(g)\) when nonempty; use an explicitly chosen fallback only when \(A(g)=\varnothing\). Align the empty-set fallback with the implementation, which uses nearest goal proximity rather than necessarily returning \(z_h\).

The official supplemental implementation and current public implementation already avoid the demonstrated failure. The latest paper formula still contains the unconditional union.

\section*{Limits}

Preliminary AI assessment; independent human scientific adjudication remains pending.\par

Analytic examples and release discrepancies do not remeasure the paper’s reported experiments.\par

\section*{Source identity}

arXiv:2509.20612v3 (latest author version confirmed from current official metadata)

\url{https://arxiv.org/pdf/2509.20612v3}

Reviewed PDF SHA-256: \nolinkurl{907db7059cb9da544e51d1401af5017181e5e71df407fbe935da6a350787301c}

The SHA-256 identifies the exact retained PDF inspected in this rapid source/argument check. Original supplied report bytes are retained locally. No unambiguous original PDF retrieval timestamp was recovered; file modification time was not substituted. Current official metadata identifies the latest author version selected here; the relevant passages were read in that version.

\end{document}
